\scanpage{18}
8) Soit maintenant $E^*$ un anneau gradué, \deleted{\unsure{muni de}} tel qu'il existe $L\in E^{(2)}$ satisfaisant aux \uline{conditions de Lefschetz} (\deleted{de} les conditions du corollaire du n° précédent). On dit que $E^*$ est un \textquotedblleft anneau \deleted{d'\unsure{opérateurs} lefschetziens \unsure{(à ne pas confondre avec la notion}} \inserted{de correspondances lefschetz\unsure{ien|iens}\textquotedblright. (Attention,} ce n'est \unsure{pourtant} qu'un \unsure{résidu} de Lefschetz), et $L$ en est un \textquotedblleft élément de Lefschetz\textquotedblright. \deleted{(\unsure{on dira simplement})\unsure{,}} L'élément $\Lambda$ associé est \deleted{appelé} \textquotedblleft l'élément \deleted{\unsure{adjoint formel}} \deleted{\unsure{anti}} \inserted{co.}lefschetz\deleted{\unsure{ien}} \textquotedblright\ associé à $L$. (NB. À montrer que $L$ se déduit de $\Lambda$ comme $\Lambda$ de $L$, par simple \unsure{symétrie|symétrique} sur les degrés de $E^*$ \deleted{$\unsure{\gtrless}-$}, et passage à l'anneau opposé --- sorte de \deleted{\illegible{1}} transposition).
% slanted
\margin{ou mieux : \textquotedblleft principe de dualité\textquotedblright\ pour les anneaux de correspondances lefschetziens.}

On peut se demander de \uline{caractériser} les \deleted{\unsure{les}} éléments de Lefschetz $L$. Une condition nécessaire, lorsque \unsure{la condition de|la condition sur} \unsure{$2n$} n'est pas \unsure{rédondant|rédondante}, i.e. que
% slanted
\margin{\unsure{désignons par|déjà par prop.|déjà pour $n$}}
\[
\begin{aligned}
\overset{\unsure{1)}}{E^{(2n)}}\ne0\quad&\deleted{\unsure{\Leftrightarrow}}\Leftrightarrow\overset{1')}{E^{(-2n)}}\ne0\deleted{\unsure{\}}}
\quad\Leftrightarrow\overset{2)}{\pi_0},\deleted{\unsure{\pi_{2n}}}\ne0\deleted{\unsure{\}}}\\
&\inserted{\Leftrightarrow\overset{2')}{\pi_{2n}}\ne0}
\Leftrightarrow\overset{3)}{L^n}\ne0\quad\Leftrightarrow\overset{3')}{\Lambda^n}\ne0\quad\deleted{\Leftrightarrow}.
\end{aligned}
\]
c'est que l'on ait $L^n\ne0$. On a parfois une réciproque, dans le cas où ($E$ étant une \uline{algèbre} sur un anneau commutatif $k$)
% archive sheet number: 17
